\documentclass[aspectratio=43,10pt]{beamer}

%% ─── Theme ────────────────────────────────────────────────────────
\usetheme[iut, sans, progressbar=segmented, nosub]{Shibui}

%% ─── Packages ────────────────────────────────────────────────────
\usepackage[utf8]{inputenc}
\usepackage[T1]{fontenc}
\usepackage[french]{babel}
\usepackage{listings}
\usepackage{translator}
\uselanguage{French}
\languagepath{French}
\usepackage{graphicx}
\usepackage{mathtools}
\usepackage{tikz}
\usepackage{pgfplots}
\pgfplotsset{compat=1.18}
\usetikzlibrary{arrows.meta,positioning,calc}
\usepgfplotslibrary{fillbetween,groupplots}

\graphicspath{{../../../common/figures/ch2/}}
\usepackage{amsmath,amsfonts,bm}

%% ─── Operators ───────────────────────────────────────────────────
\DeclareMathOperator*{\argmax}{arg\,max}
\DeclareMathOperator*{\argmin}{arg\,min}
\DeclareMathOperator{\MSE}{MSE}
\DeclareMathOperator{\sign}{sign}
\DeclareMathOperator{\sinc}{sinc}
\DeclareMathOperator{\tr}{tr}

%% ─── Number sets ─────────────────────────────────────────────────
\def\R{{\mathbb{R}}}
\def\C{{\mathbb{C}}}
\def\N{{\mathbb{N}}}
\def\Z{{\mathbb{Z}}}

%% ─── Infinitesimal and units ────────────────────────────────────
\def\d{{\rm d}}
\def\Hz{{\mbox{Hz}}}
\def\m{{\mbox{m}}}

%% ─── Complex / signal ───────────────────────────────────────────
\DeclareMathOperator{\complexj}{\mathsf{j}}   % imaginary unit

%% ─── Probability ─────────────────────────────────────────────────
\newcommand{\Esp}[2]{\mathbb{E}_{#1}\!\left[#2\right]}
\renewcommand{\Pr}{\mathbb{P}}

%% ─── Linear algebra ──────────────────────────────────────────────
\newcommand{\abs}[1]{\left|{#1}\right|}
\newcommand{\norm}[1]{\left\|#1\right\|}
\newcommand{\inner}[2]{\left\langle#1,\,#2\right\rangle}

%% ─── Logic ───────────────────────────────────────────────────────
\newcommand{\then}{\Rightarrow}
\newcommand{\iif}{\Leftrightarrow}

%% ─── Draft annotation ────────────────────────────────────────────
\newcommand{\red}[1]{{\color{red}#1}}

%% ─── Shared figure color (used by common/figures/*/*.tex bodies) ──
\definecolor{niceblue}{RGB}{0,128,172}

%% ─── Rappel header (styled subtitle for reference-card asides) ───
%% Requires \definecolor{niceblue}{...} to be defined by the including document.
\newcommand{\rappelhead}[1]{%
    \bigskip
    {\color{niceblue}\bfseries #1}\\[-0.4em]
    {\color{niceblue!35}\rule{\linewidth}{0.6pt}}
    \smallskip

}


%% ─── Metadata ─────────────────────────────────────────────────────
\title{Outils Mathématiques et Logiciel --- OML3}
\subtitle{Chapitre 2 --- Transformée de Fourier}
\author{\textbf{BUT2} -- Génie Électrique et Informatique Industrielle}
\date{2025-2026}

%% ─── Theorem environments ─────────────────────────────────────────
%% `definition` and `theorem` are already provided (and auto-localized
%% via babel) by the beamer class itself -- only declare the extra ones.
\newtheorem{proposition}{Propriété}
\newtheorem{remark}{Remarque}

%% ─── Text helper (mirrors poly/main.tex) ───────────────────────────
\newcommand{\highlight}[1]{{\color{niceblue}\emph{#1}}}
\newcommand{\tf}{\overset{\mathcal F}{\longleftrightarrow}}

%% ─── Figure helper : fit any \input'd figure to a target width ────
\newcommand{\slidefig}[2]{%
\begin{center}
\resizebox{#1}{!}{\input{../../../common/figures/ch2/#2}}
\end{center}
}

%% ─────────────────────────────────────────────────────────────────
\begin{document}
\shorthandoff{:;!?}

\maketitle

%% ═══════════════════════════════════════════════════════════════════
%% SÉANCE 1 (CM1, 1h) : sections "Introduction et contexte",
%%                       "Transformée de Fourier"
%% SÉANCE 2 (CM2, 1h) : sections "Impulsion de Dirac",
%%                       "Produit de convolution"
%% Mirrors poly/main.tex. Exemple fil rouge : le signal porte
%% Pi_theta, dont la TF (theta sinc(pi f theta)) sert d'illustration a
%% plusieurs propriétés (dualité, échelle, parité, Parseval) plutôt que
%% de multiplier les exemples.
%% ═══════════════════════════════════════════════════════════════════

%% ─── SECTION : Introduction et contexte ─────────────────────────
\section{Introduction et contexte}

\begin{frame}{Et les signaux non périodiques ?}
    \begin{block}{Chapitre 1}
        Séries de Fourier : décomposer un signal \highlight{périodique}
        en somme de sinusoïdes.
    \end{block}
    \begin{block}{Chapitre 2}
        Peut-on faire de même avec un signal \highlight{non périodique}  ?
        \slidefig{0.6\textwidth}{context-signal-example.tex}
        \begin{alertblock}{}
        \end{alertblock}
    \end{block}
\end{frame}

\begin{frame}{Intégrale sur $\R$ et sinus cardinal}
    \begin{definition}[Intégrale sur $\R$]
        $\displaystyle\int_{-\infty}^{+\infty} g(t)\,\d t = \lim_{\substack{A\to+\infty\\B\to+\infty}} \int_{-A}^{B} g(t)\,\d t$,
        lorsque cette limite existe.
    \end{definition}
    \begin{definition}[Sinus cardinal]
        $\sinc(x) = \dfrac{\sin(x)}{x}$ si $x\neq0$, et $\sinc(0)=1$.
    \end{definition}
    \slidefig{0.55\textwidth}{sinc-def.tex}
\end{frame}

\begin{frame}{Énergie et puissance d'un signal}
    \begin{definition}[Énergie, puissance moyenne]
        $$
            E_x = \int_{-\infty}^{+\infty} |x(t)|^2\,\d t, \qquad
            P_x = \lim_{T\to+\infty} \frac1T\int_{-T/2}^{T/2} |x(t)|^2\,\d t.
        $$
    \end{definition}
    \slidefig{0.85\textwidth}{energy-vs-power.tex}
    \begin{center}
        \small Signal périodique : $P_x$ finie, $E_x$ infinie. \quad
        Signal à support fini : $E_x$ finie.
    \end{center}
\end{frame}

\begin{frame}{Un exemple particulier}
    % \begin{center}
    %     Un signal à énergie finie comme \highlight{limite} d'un signal
    %     périodique dont la période $T\to+\infty$ : un train
    %     d'impulsions de période $T$, largeur $\theta$, hauteur
    %     $1/\theta$.
    % \end{center}
    \begin{overprint}
    \onslide<1>
    \slidefig{\textwidth}{intro-limit-small-T.tex}
    \begin{center}\small Fonction périodique $\to$   Coefficients complexe de sa série de Fourier\end{center}
    \onslide<2>
    \slidefig{\textwidth}{intro-limit-large-T.tex}
    \begin{center}\small Lorsque la période $T$ augmente les raies $Tc_n$ se resserrent et
    dessinent une enveloppe continue.\end{center}
    \end{overprint}
\end{frame}



%% ─── SECTION : Transformée de Fourier ───────────────────────────
\section{Transformée de Fourier}

\begin{frame}{Définition : TF et TFI}
    \begin{definition}[Transformée de Fourier]
        $$
            X(f) = \mathcal F[x](f) = \int_{-\infty}^{+\infty} x(t)\,e^{-2\complexj\pi ft}\,\d t
        $$
    \end{definition}
    \begin{definition}[Transformée de Fourier inverse]
        $$
            x(t) = \mathcal F^{-1}[X](t) = \int_{-\infty}^{+\infty} X(f)\,e^{2\complexj\pi ft}\,\d f
        $$
    \end{definition}
    \begin{center}
        \small Seul le signe de l'exposant change entre les deux.
        \quad $x(t)\tf X(f)$
    \end{center}
\end{frame}

\begin{frame}{De la série à l'intégrale}
    \[
        x(t) = \sum_{n=-\infty}^{+\infty} c_n\,e^{\complexj2\pi nf_0 t}
        = \sum_{n=-\infty}^{+\infty} X_T(nf_0)\,e^{\complexj2\pi nf_0 t}\ f_0,
    \]
    $$
     X_T(nf_0) f_0 = c_n
    $$
    \slidefig{0.68\textwidth}{riemann-sum.tex}
    \begin{center}
        \small Quand $T\to+\infty$, $f_0\to\d f$ : la somme des raies
        devient l'intégrale
    \end{center}
\end{frame}

\begin{frame}{Exemple : le signal porte}
    \begin{definition}[Signal porte]
        $$
            \Pi_\theta(t) = \begin{cases} 1 & \text{si } |t|<\theta/2 \\ 0 & \text{si } |t|\geq\theta/2 \end{cases}
        $$
    \end{definition}
    \slidefig{\textwidth}{porte-sinc.tex}
    \begin{center}
        $\Pi_\theta(t) \tf \theta\,\sinc(\pi f\theta)$ 
    \end{center}
\end{frame}

\begin{frame}{Deux domaines, une seule information}
    \begin{center}
        \Large $\mathcal F^{-1}\big[\mathcal F[x]\big] = x$
    \end{center}
    \vspace{1em}
    \begin{center}
        \small Domaine \highlight{temporel} ($x(t)$) et domaine
        \highlight{fréquentiel} ($X(f)$) 
        
        Deux représentations
        équivalentes du même signal.
    \end{center}
    % \begin{alertblock}{Condition d'existence (simplifiée)}
    %     Si $x$ est intégrable ($\int|x(t)|\,\d t<\infty$), $X(f)$
    %     existe pour tout $f$. Cela exclut $\cos$, $\sin$, une
    %     constante --- on lèvera cette limitation en \S3 avec
    %     l'impulsion de Dirac.
    % \end{alertblock}
\end{frame}

\begin{frame}{Propriété : dualité}
    \begin{proposition}[Dualité]
        Si $x(t)\tf X(f)$, alors $\mathcal F[X](f) = x(-f)$.
    \end{proposition}
    \slidefig{0.85\textwidth}{dualite-2x2.tex}
    \begin{center}
        \small En haut : $\Pi_\theta(t)\tf\theta\sinc(\pi f\theta)$. En
        bas : le même sinus cardinal, relu comme un signal
        \highlight{temporel} (variable $t$), a pour TF une porte, cette
        fois en \highlight{fréquence} (variable $f$) --- $\Pi_\theta$
        étant paire, $x(-f)=x(f)$.
    \end{center}
\end{frame}

\begin{frame}{Propriétés : linéarité, décalages}
    \begin{proposition}[Linéarité]
        $\alpha x(t) + \beta y(t) \tf \alpha X(f) + \beta Y(f)$
    \end{proposition}
    \begin{proposition}[Décalage temporel et fréquentiel]
        $$
            x(t-t_0) \tf X(f)\,e^{-2\complexj\pi ft_0}, \qquad\qquad
            x(t)\,e^{2\complexj\pi f_0t} \tf X(f-f_0)
        $$
    \end{proposition}
    \begin{center}
        \small Un décalage temporel ne change pas $|X(f)|$, seulement
        sa phase.
    \end{center}
\end{frame}

\begin{frame}{Propriété : changement d'échelle}
    \begin{proposition}[Changement d'échelle]
        $$
            x(at) \tf \frac{1}{|a|}X\!\left(\frac fa\right)
        $$
    \end{proposition}
    \slidefig{0.7\textwidth}{echelle-small-theta.tex}
    \slidefig{0.7\textwidth}{echelle-large-theta.tex}
    \begin{center}
        \small
        Comprimer dans le temps dilate le spectre, et réciproquement.
    \end{center}
\end{frame}

\begin{frame}{Propriété : dérivation}
    \begin{proposition}[Dérivation]
        Si $x(t)\to0$ en $\pm\infty$ et $x(t)\tf X(f)$, alors
        $$
            x'(t) \tf 2\complexj\pi f\,X(f).
        $$
    \end{proposition}
    \begin{center}
        \small Dériver dans le temps $\iff$ multiplier par
        $2\complexj\pi f$ en fréquence
        
        Intégrer $\iff$ diviser (à
        une constante près).
    \end{center}
\end{frame}

\begin{frame}{Propriété : parité d'un signal réel}
    \begin{proposition}[Parité et TF d'un signal réel]
        Soit $x$ une fonction réel, $x(t)\tf X(f)$.
        \begin{itemize}
            \item $X(-f) = \overline{X(f)}$, 
            \item $x$ paire $\Rightarrow$ $X$ réelle et paire, 
            \item $x$ impaire $\Rightarrow$ $X$ imaginaire pure et impaire.
        \end{itemize}
    \end{proposition}

\end{frame}

\begin{frame}{Propriété : Parseval-Plancherel}
    \begin{proposition}[Parseval-Plancherel]
        $$
            E_x = \int_{-\infty}^{+\infty} |x(t)|^2\,\d t = \int_{-\infty}^{+\infty} |X(f)|^2\,\d f
        $$
    \end{proposition}
    \slidefig{\textwidth}{parseval-porte.tex}
    \begin{center}
        \small Deux courbes de formes très différentes, même aire
        $E_x=\theta$ : $|X(f)|^2$ est une densité spectrale d'énergie.
    \end{center}
\end{frame}

%% ─── SECTION : Impulsion de Dirac ────────────────────────────────
\section{Impulsion de Dirac}

\begin{frame}{L'impulsion de Dirac}
    \begin{definition}[Impulsion de Dirac]
        Pour $\theta>0$, soit $p_\theta = \frac1\theta\Pi_\theta$ (aire
        $1$, largeur $\theta$). L'impulsion de Dirac est
        $$
            \delta(t) = \lim_{\theta\to0} p_\theta(t).
        $$
    \end{definition}
    \begin{overprint}
    \onslide<1>
    \slidefig{\textwidth}{dirac-as-limit.tex}
    \begin{center}\small $p_\theta$ se resserre et grandit, mais garde
    toujours une aire égale à $1$.\end{center}
    \onslide<2>
    \slidefig{0.4\textwidth}{dirac-arrow-single.tex}
    \begin{center}\small Convention graphique : une flèche verticale

    (\red{Attention} : $\delta$ n'est pas une fonction usuelle).\end{center}
    \end{overprint}
\end{frame}

\begin{frame}{L'outil central : propriété d'échantillonnage}
    \begin{proposition}[Propriétés de $\delta$]
        Pour $x$ continue en $t_0$ :
        $$
            x(t)\,\delta(t-t_0) = x(t_0)\,\delta(t-t_0), \qquad
            \int_{-\infty}^{+\infty} x(t)\,\delta(t-t_0)\,\d t = x(t_0).
        $$
    \end{proposition}
    \slidefig{0.4\textwidth}{dirac-sampling.tex}
    \begin{alertblock}{Retenir}
        $\delta$ « prélève » la valeur de $x$ en un point --- c'est
        cette propriété qui rend $\delta$ si utile.
    \end{alertblock}
\end{frame}

\begin{frame}{Changement d'échelle et parité de $\delta$}
    \begin{proposition}[Changement d'échelle]
        $$
            \delta(at) = \frac{1}{|a|}\delta(t), \qquad\qquad \delta(-t)=\delta(t) \ \text{($\delta$ est paire)}.
        $$
    \end{proposition}
\end{frame}

\begin{frame}{TF d'une impulsion}
    \begin{proposition}[TF d'une impulsion]
        $$
            \mathcal F[\delta](f) = 1, \qquad\qquad \mathcal F[1](f) = \delta(f).
        $$
    \end{proposition}
    \begin{overprint}
    \onslide<1>
    \slidefig{\textwidth}{dirac-tf-limit.tex}
    % \begin{center}\small Intuition : $\mathcal F[p_\theta](f)=\sinc(\pi f\theta)$
    % (gris $\theta=0.5$, bleu $\theta=0.15$) s'aplatit vers la constante
    % $1$ (rouge) quand $\theta\to0$.\end{center}
    \onslide<2>
    \slidefig{\textwidth}{dirac-tf-constant.tex}
    \begin{center}\small Une impulsion contient, avec le même poids,
    toutes les fréquences.\end{center}
    \end{overprint}
\end{frame}

\begin{frame}{Exemple : TF d'une impulsion décalée}
    \[
        \mathcal F[\delta(\cdot-t_0)](f) = e^{-2\complexj\pi ft_0}
    \]
    \slidefig{\textwidth}{dirac-tf-decale.tex}

    % \begin{center}
    %     \small Module constant égal à $1$ : c'est la propriété de
    %     \highlight{décalage temporel} du \S2 --- décaler ne change que
    %     la phase.
    % \end{center}
\end{frame}

\begin{frame}{Exemple : TF de $\cos$ et $\sin$}
    \begin{columns}
        \column{0.5\textwidth}
        \small
        \textbf{Chapitre 1 (rappel)}
        \[
            \cos(\omega_0t) = \tfrac12e^{\complexj\omega_0t} + \tfrac12e^{-\complexj\omega_0t}
        \]
        Spectre \highlight{discret} : deux raies en $\pm f_0$, poids
        $c_1=c_{-1}=\frac12$.
        \column{0.5\textwidth}
        \small
        \textbf{Chapitre 2 (ici)}
        \[
            \mathcal F[\cos(2\pi f_0t)](f) = \tfrac12\delta(f-f_0) + \tfrac12\delta(f+f_0)
        \]
        Spectre \highlight{continu} : mêmes fréquences, mêmes poids,
        mais des impulsions.
    \end{columns}
    \slidefig{0.75\textwidth}{cos-sin-tf-spectrum.tex}
\end{frame}

%% ─── SECTION : Produit de convolution ────────────────────────────
\section{Produit de convolution}

\begin{frame}{Le produit de convolution}
    \begin{definition}[Produit de convolution]
        $$
            h(t) = (x*y)(t) = \int_{-\infty}^{+\infty} x(u)\,y(t-u)\,\d u
        $$
    \end{definition}
    \begin{columns}
        \column{0.5\textwidth}
        \slidefig{\textwidth}{conv-setup.tex}
        % \begin{center}\small $x$ et $y$, avant tout décalage.\end{center}
        \column{0.5\textwidth}
        \slidefig{\textwidth}{conv-product.tex}
        % \begin{center}\small Pour $t=a>0$ : $x(u)\,y(a-u)$, aire
        % colorée $=(x*y)(a)$.\end{center}
    \end{columns}
    \begin{center}
        \small Bilinéaire et commutatif : $x*y=y*x$.
    \end{center}
\end{frame}

\begin{frame}{Invariance temporelle}
    \begin{overprint}
        \onslide<1>
        \begin{proposition}[Linéarité]
            $\forall t,t_0\in\R$, $$((x_1 + x_2)*y)(t) = (x_{1} * y)(t) + (x_{2} * y)(t).$$
        \end{proposition}
        \slidefig{0.85\textwidth}{lti-linearity.tex}
        \begin{center}\small Linéarité de $S$.\end{center}
        \onslide<2>
        \begin{proposition}[Invariance temporelle]
            $\forall t,t_0\in\R$, $$(x*y)(t-t_0) = (x_{-t_{0}} * y)(t),$$ avec $x_{-t_{0}}(u) = x(u-t_0)$.
        \end{proposition}
    \slidefig{0.7\textwidth}{lti-time-invariance.tex}
    \begin{center}\small Invariance temporelle de $S$.\end{center}
    \end{overprint}
    \begin{alertblock}{Pourquoi ça compte}
        La sortie d'un système \highlight{linéaire invariant dans le
        temps} (un filtre) s'obtient par convolution de l'entrée $x$
        avec sa \highlight{réponse impulsionnelle} $h$ : $y=x*h$.
    \end{alertblock}
\end{frame}

\begin{frame}{Exemples de systèmes LTI}
    \begin{columns}
        \column{0.5\textwidth}
        \centering \small Circuit RC (filtre)
        \slidefig{\textwidth}{example-rc.tex}
        \begin{center}\small $v_e(t) \ \longrightarrow\ v_s(t)$\end{center}
        \column{0.5\textwidth}
        \centering \small Moteur à courant continu
        \slidefig{\textwidth}{example-motor.tex}
        \begin{center}\small $V(t) \ \longrightarrow\ \omega(t)$\end{center}
    \end{columns}
    \begin{center}
        \small Deux systèmes physiques très différents, mais tous deux
        \highlight{linéaires} et \highlight{invariants dans le temps}.
    \end{center}
\end{frame}

\begin{frame}{Convolution avec une porte, avec une impulsion}
    \begin{proposition}[Convolution avec porte et impulsion]
        \begin{itemize}
            \item $x*\delta = x$ : $\delta$ est neutre pour $*$ ;
            \item $x*\delta(\cdot-t_0) = x(\cdot-t_0)$ : décale $x$ ;
            \item $x*\Pi_\theta$ lisse $x$ (moyenne glissante de largeur $\theta$).
        \end{itemize}
    \end{proposition}
    \slidefig{.9\textwidth}{conv-shift-illustration.tex}
\end{frame}

\begin{frame}{Dualité produit / convolution}
    \begin{proposition}[Dualité produit / convolution]
        $$
            (x*y)(t) \tf X(f)\,Y(f), \qquad\qquad x(t)\,y(t) \tf (X*Y)(f)
        $$
    \end{proposition}
    \begin{center}
        Une convolution dans un domaine devient un simple
        \highlight{produit} dans l'autre.
    \end{center}
\end{frame}

\begin{frame}{Pour finir : observer un cosinus pendant une durée finie}
    \[
        \cos(2\pi f_0t)\times\Pi_\theta(t) \quad\longrightarrow\quad
        \frac\theta2\Big(\sinc(\pi\theta(f-f_0)) + \sinc(\pi\theta(f+f_0))\Big)
    \]
    \slidefig{\textwidth}{windowed-cosine-spectrum.tex}
    \begin{center}
        \small Produit dans le temps = convolution en fréquence : les
        impulsions idéales de $\cos$ s'élargissent en lobes de sinus
        cardinal. Observer un signal pendant un temps fini élargit
        toujours ses raies spectrales.
    \end{center}
\end{frame}

\end{document}
